% Discovery flow in scripts/run_bertopic.py. Existing vectors are reused;
% cleaned text supplies the c-TF-IDF descriptions. The input example reuses the
% verified saved page and vector from Figure 5; the flow describes the full sample.
\noindent\begin{minipage}{\linewidth}
\centering
\begin{tikzpicture}[
  stage/.style = {rounded corners=4pt, draw, line width=0.8pt,
                  text width=2.50cm, minimum height=2.30cm,
                  align=center, inner xsep=4pt, inner ysep=5pt,
                  font=\small\sffamily},
  flow/.style = {-{Stealth[length=2.2mm,width=1.7mm]},
                 line width=0.9pt, draw=black!65},
  side/.style = {flow, dashed},
]

\node[font=\small\sffamily\bfseries, text=accent, align=center,
      text width=6.2cm] at (3.2,3.65)
  {Run separately for\\English and German};

\node[stage, minimum height=4.50cm, fill=fillblue, draw=boxblue] (input) at (0,0)
  {\textbf{1 Input}\\[4pt]
   Saved E5 vectors\\
   One real example:\\[5pt]
   {\footnotesize\ttfamily
    [\\
    0.008740,\\
    0.031677,\\
    -0.000636,\\
    -0.018505,\\
    0.019698,\\
    \ldots ]}};

\node[stage, fill=fillgreen, draw=boxgreen] (reduce) at (3.2,0)
  {\textbf{2 UMAP}\\[4pt]
   Reduce the\\number of\\coordinates};

\node[stage, fill=fillgreen, draw=boxgreen] (cluster) at (6.4,0)
  {\textbf{3 HDBSCAN}\\[4pt]
   Find dense\\groups of\\nearby vectors};

\node[stage, fill=black!3, draw=diagray] (words) at (9.6,0)
  {\textbf{4 Keywords}\\[4pt]
   Count words\\by topic; rank\\with c-TF-IDF};

\node[stage, minimum height=4.50cm, fill=fillblue, draw=boxblue] (output) at (12.8,0)
  {\textbf{5 Output}\\[4pt]
   {\footnotesize\textbf{Shown page}\\[3pt]
    \ttfamily
    uuid: ...3b2f\\
    topic\_id: -1\\
    topic\_weights:\\
    0.0}\\[6pt]
   {\footnotesize\textbf{Another topic: 85}\\[3pt]
    robot, path,\\
    control, \ldots}};

\draw[flow] (input.east) -- (reduce.west);
\draw[flow] (reduce.east) -- (cluster.west);
\draw[flow] (cluster.east) -- (words.west);
\draw[flow] (words.east) -- (output.west);

\node[stage, text width=4.70cm, minimum width=5.10cm,
      minimum height=2.60cm, fill=black!3, draw=diagray,
      align=left, font=\footnotesize\sffamily] (text) at (9.6,4.10)
  {\textbf{Cleaned body text}\\[4pt]
   Example from the same page:\\[3pt]
   \textit{This section shows video clips of experiments in which the path of robot arm
   is refined by camera images.}};
\draw[side, draw=diagray] (text.south) -- (words.north);

\node[stage, text width=7.20cm, minimum height=1.15cm,
      fill=fillred, draw=boxred] (noise) at (6.4,-2.55)
  {\textbf{Noise label $-1$}\\[3pt]
   Keep these pages for labeling (Section~\ref{sec:method-coverage})};
\draw[side, draw=boxred] (cluster.south) -- (noise.north);

\end{tikzpicture}
\captionof{figure}{BERTopic discovery on one language sample. \measured{} The input examples
reuse the robotics page from Figure~\ref{fig:embedding-chunks}: its saved vector and a
cleaned-text excerpt. Only the first five vector values are shown, rounded; the vector
represents the full page. All sampled pages are processed together. The displayed page
received $-1$ during discovery and continues to document labeling
(Section~\ref{sec:method-coverage}). Its output row comes from \texttt{doc\_topics.parquet},
with the page ID shortened here. The separate keyword example for topic 85 comes from
\texttt{topic\_top\_words.json} and describes other sampled pages.}
\label{fig:bertopic-pipeline}
\end{minipage}
